% =============================================================================
%  5-hipotesis.tex — Balance de la hipótesis
% =============================================================================

\section{Confirmación Parcial de la Hipótesis}

\begin{frame}{Balance de la hipótesis}
  \begin{columns}[T,onlytextwidth]
    \begin{column}{0.48\textwidth}
      \begin{block}{Se confirma}
        \begin{itemize}
          \item La red radial-concéntrica concentra los caminos más cortos en
                pocos nodos centrales.
          \item Con el Anillo Metro (Caso~3), el tiempo medio de viaje modelado
                ($T$) baja de $113.91$ a $101.92$~min.
          \item La intermediación de Tacubaya cae
                $29.7\,\text{\%}$ frente a la red base.
        \end{itemize}
      \end{block}
    \end{column}
    \begin{column}{0.48\textwidth}
      \begin{alertblock}{Se matiza}
        \begin{itemize}
          \item La cobertura metropolitana ($C$) pasa de
                $13.29\,\text{\%}$ a $13.43\,\text{\%}$.
                Aún quedan 41 municipios sin estación a distancia caminable.
          \item El Gini de intermediación sube de $0.8007$ a $0.8134$.
                La concentración se desplaza hacia nodos del anillo.
        \end{itemize}
      \end{alertblock}
    \end{column}
  \end{columns}

  \medskip
  {\footnotesize\itshape Fuente: resultados de la tesis, caps.~6 y~7.\par}
\end{frame}

\begin{frame}{Fortalezas y hallazgos}
  \begin{enumerate}
    \item \textbf{Soberanía tecnológica.}
          Apímetro integra los datos, VFTModel calcula los indicadores y
          Transport-gis permite visualizarlos. El cálculo puede reproducirse
          sin licencias comerciales.

    \item \textbf{Programa Maestro del Metro de 1985.}
          El programa descartó los dos corredores del Periférico por su baja
          captación de demanda. Entre los criterios figuraba la captación por
          kilómetro. El modelo actual muestra cómo el Anillo Metro reorganiza
          la conectividad metropolitana.

    \item \textbf{CETRAM de nueva generación.}
          Seis estaciones del Anillo Metro alcanzan una intermediación
          $B(v)$ de $0.12$ a $0.16$ con 6 a 12 conexiones físicas.
          Conviene planearlas como nodos intermodales, con rutas alimentadoras
          y accesos de última milla.
  \end{enumerate}

  \medskip
  {\footnotesize\itshape Fuentes: Programa Maestro del Metro (DDF/COVITUR, 1985)
    y resultados de la tesis, caps.~6 y~7.\par}
\end{frame}

\begin{frame}{Limitaciones del modelo}
  \begin{alertblock}{Alcance de los resultados}
    \begin{itemize}
      \item \textbf{Enfoque topológico.}
            El modelo estima recorridos y tiempos de viaje. Quedan fuera las
            matrices origen--destino, la afluencia y la densidad poblacional.

      \item \textbf{Factor de desviación ($DI$).}
            La red base tiene una mediana de $1.52$. La comparación con los
            Casos~2 y~3 quedó pendiente porque las muestras de pares
            origen--destino eran distintas.

      \item \textbf{Vulnerabilidad estructural ($\Delta E$).}
            Se exploró retirando, uno por uno, cinco nodos críticos del grafo
            base. El cálculo quedó pendiente en la capa final del SIG.
    \end{itemize}
  \end{alertblock}

  \medskip
  {\footnotesize\itshape Fuente: metodología y resultados de la tesis,
    caps.~5 a~7.\par}
\end{frame}
